\documentclass[UTF8,zihao=-4]{ctexart}

\usepackage[a4paper,margin=2.4cm]{geometry}
\usepackage{array}
\usepackage{booktabs}
\usepackage{longtable}
\usepackage{tabularx}
\usepackage{enumitem}
\usepackage{xcolor}
\usepackage{listings}
\usepackage{graphicx}
\usepackage{hyperref}
\usepackage{fancyhdr}
\usepackage{titlesec}
\usepackage{float}

\definecolor{Primary}{HTML}{174A7C}
\definecolor{Accent}{HTML}{2E7D6E}
\definecolor{Soft}{HTML}{F2F6FA}
\definecolor{Warn}{HTML}{8A5A00}

\hypersetup{
  unicode=true,
  colorlinks=true,
  linkcolor=Primary,
  urlcolor=Accent,
  pdftitle={信息学竞赛 AI 自动造数据与数据校验工具项目设计报告},
  pdfauthor={专业实习项目组}
}

\pagestyle{fancy}
\setlength{\headheight}{15pt}
\fancyhf{}
\lhead{信息学竞赛 AI 自动造数据与数据校验工具}
\rhead{项目设计报告}
\cfoot{\thepage}

\titleformat{\section}{\Large\bfseries\color{Primary}}{\thesection}{0.8em}{}
\titleformat{\subsection}{\large\bfseries\color{Accent}}{\thesubsection}{0.8em}{}
\titleformat{\subsubsection}{\normalsize\bfseries}{\thesubsubsection}{0.8em}{}

\setlist[itemize]{leftmargin=2em,itemsep=0.2em,topsep=0.3em}
\setlist[enumerate]{leftmargin=2em,itemsep=0.2em,topsep=0.3em}
\renewcommand{\arraystretch}{1.28}
\setlength{\parskip}{0.45em}
\setlength{\parindent}{2em}
\emergencystretch=2em

\lstdefinestyle{reportcode}{
  basicstyle=\ttfamily\small,
  backgroundcolor=\color{Soft},
  frame=single,
  breaklines=true,
  columns=fullflexible,
  keepspaces=true,
  showstringspaces=false,
  tabsize=2
}
\lstset{style=reportcode}

\newcommand{\docversion}{V1.0}
\newcommand{\projectname}{信息学竞赛 AI 自动造数据与数据校验工具}

\begin{document}

\begin{titlepage}
  \centering
  \vspace*{2.2cm}
  {\Huge\bfseries \projectname \par}
  \vspace{0.8cm}
  {\LARGE 完整项目设计报告\par}
  \vspace{1.2cm}
  {\large 专业实习小学期项目设计阶段\par}
  \vspace{0.5cm}
  {\large 版本：\docversion\par}
  \vfill
  \begin{tabular}{rl}
    项目类型： & 中学 / 信息学竞赛工具 \\
    设计目标： & 指导 MVP 开发、答辩演示与需求二次确认 \\
    编制日期： & 2026 年 7 月 9 日 \\
  \end{tabular}
  \vfill
\end{titlepage}

\tableofcontents
\newpage

\section{项目背景与建设目标}

专业实习课程要求项目组完成组队选题、项目设计、项目开发、项目提交与项目答辩。课程考查文件明确指出，项目设计部分占总成绩 20\%，设计报告需要体现项目业务逻辑结构清晰、功能完善、界面设计合理美观，并能够反馈需求方进行二次需求确认。当前报告面向项目设计阶段，目标是把需求、业务流程、系统架构、模块职责、界面方案、MVP 边界和验收标准沉淀为可执行的开发依据。

本项目的题目来源为“信息学竞赛 AI 自动造数据与数据校验工具”。原始需求聚焦中小学信息学竞赛训练场景，痛点包括人工造数据效率低、边界数据覆盖不全、测试数据容易出错、标准输出生成繁琐以及数据校验不足。项目目标是借助 AI 分析与规划能力，在本地实现一个 Polygon-like 的单题测试数据生成与验题流程。

本项目不是直接对接在线评测系统，也不是完整复制 Codeforces Polygon。MVP 阶段只围绕单道题完成数据生成、验证、标准输出生成、质量报告与本地导出。系统默认输入不是完整题面，而是用户导入的 C++ 标准程序、数据范围、数量规模和难度等级。AI 负责分析、规划、生成代码草稿和修复建议；Docker、testlib validator、jngen generator 和用户标程负责确定性执行与校验。

\subsection{建设目标}

\begin{enumerate}
  \item 建立清晰的单题数据生成业务流程，覆盖输入配置、标程导入、数据结构识别、测试点计划、生成器与校验器生成、Docker 执行、报告与导出。
  \item 建立安全可控的执行机制，所有不可信 C++ 代码必须在 Docker 沙箱内编译和运行，禁止在宿主机直接执行。
  \item 建立以 testlib validator 为核心的数据合法性校验机制，确保最终数据包中的每个 \texttt{.in} 文件都通过严格验证。
  \item 建立以用户导入标程为唯一答案来源的标准输出生成机制，不在 MVP 中由 AI 自动生成标准程序。
  \item 建立面向需求方的二次确认闭环，把不确定输入结构、数据范围表达方式、OJ 兼容、课堂练习、模拟比赛等问题单独列出。
  \item 建立可指导后续 MVP 开发的技术方案，包括 FastAPI、React + Vite、SQLite、Docker、Dify Python client wrapper、testlib.h、jngen.h、本地文件存储和 zip 导出。
\end{enumerate}

\section{用户角色与使用场景}

\subsection{用户角色}

\begin{longtable}{p{3.3cm}p{5.3cm}p{6.2cm}}
\toprule
角色 & 主要目标 & 关注点 \\
\midrule
信息学教师 / 出题人 & 为单道训练题快速生成可靠测试数据 & 数据边界覆盖、格式正确、输出可信、导出方便 \\
竞赛训练负责人 & 批量准备训练材料并检查数据质量 & 流程可控、报告可读、失败原因明确 \\
项目组开发人员 & 根据设计报告开发 MVP & 模块边界、接口定义、数据模型、执行安全 \\
需求方 / 学校联系人 & 判断方案是否符合教学训练场景 & 业务逻辑、功能完整度、界面是否易用、后续需求边界 \\
\bottomrule
\end{longtable}

\subsection{典型使用场景}

\begin{enumerate}
  \item 教师创建单题项目，填写项目名称和说明。
  \item 教师导入 C++ 标准程序，系统在 Docker 中编译，确认标程可执行。
  \item 教师填写数据范围文本、数量规模数值，并从单选下拉框中选择难度等级。数据范围文本可采用类似 Constraints 的描述，变量名需与标程代码保持一致。
  \item 系统调用 Dify / Mock Dify 工作流，分析标程输入读取逻辑并生成输入结构草案、数据结构识别结果和测试点计划。
  \item 系统生成 generator.cpp 和 validator.cpp 草稿，用户可预览并人工确认。
  \item 系统在 Docker 中编译 generator、validator 和 solution，批量生成 \texttt{.in} 文件。
  \item 系统使用 testlib validator 严格校验每个输入，过滤不合法、重复或超限测试点。
  \item 系统使用用户标程生成 \texttt{.out} 文件，并记录运行时间、内存、退出码和文件大小。
  \item 系统生成质量报告，展示数据结构覆盖、边界覆盖、运行日志和需要人工确认的问题。
  \item 用户导出本地 zip 数据包，用于后续教学、训练或进一步人工整理。
\end{enumerate}

\section{需求分析}

\subsection{功能需求}

\begin{longtable}{p{3.5cm}p{10.5cm}}
\toprule
功能模块 & 功能说明 \\
\midrule
单题项目管理 & 创建、查看、编辑单题项目，保存项目状态和基础配置。 \\
标程导入与编译 & 上传或粘贴 C++ 标准程序，使用 Docker 编译，记录编译结果和错误日志。 \\
输入配置管理 & 通过多行文本框录入类似 Constraints 的数据范围说明，变量名需与标程代码保持一致；通过单行文本框录入测试点数量等数值规模；通过单选下拉框选择难度等级；不要求用户输入完整题面。 \\
Dify 工作流 & 分析标程读取逻辑，结合输入配置输出输入结构草案、数据结构识别和测试点计划。 \\
测试点计划 & 每个测试点包含 category、scale、purpose、seed、resource limit。 \\
generator 生成 & 生成基于 jngen 的 generator.cpp 草稿，支持 seed 和 case\_id。 \\
validator 生成 & 生成基于 testlib 的 validator.cpp 草稿，严格检查输入格式、范围和结构约束。 \\
Docker 编译运行 & 编译和运行 generator、validator、solution，限制 CPU、内存、时间和网络。 \\
输入校验 & 对每个 \texttt{.in} 执行文件大小检查、基础格式检查、testlib validator 校验、重复检测和资源限制检查。 \\
答案生成 & 使用用户标程生成 \texttt{.out}，记录 exit code、runtime、memory、stdout size、stderr、超时和内存超限状态。 \\
质量报告 & 生成 report.md / 页面报告，展示配置摘要、测试点计划、验证结果、执行资源和失败原因。 \\
zip 导出 & 导出本地通用数据包，包括输入输出文件、generator.cpp、validator.cpp、solution.cpp、metadata.json 和 report.md。 \\
\bottomrule
\end{longtable}

\subsection{非功能需求}

\begin{itemize}
  \item 安全性：不可信代码必须在 Docker 中编译和运行，禁止网络访问，只挂载临时工作目录。
  \item 可复现性：generator 必须支持 seed 和 case\_id，测试点生成过程应可追踪。
  \item 可观测性：Dify 调用、Docker 编译、运行日志、validator 失败原因和标程失败原因都必须记录。
  \item 可维护性：后端使用 FastAPI 分层设计，前端使用 React + Vite，数据持久化使用 SQLite。
  \item 可解释性：AI 推断出的输入结构和隐含约束必须展示不确定字段，不能直接替代人工确认。
  \item 性能边界：系统必须配置时间、内存、文件大小、总数据包大小和最大重试次数。
\end{itemize}

\subsection{需求边界说明}

当前应用优先覆盖单题数据生成、输入校验、标准输出生成、质量报告和 zip 导出闭环。暂不纳入应用且需要二次确认的事项集中列于“应用功能边界”章节，避免在需求分析阶段过早承诺 OJ 兼容、课堂练习、模拟比赛、SPJ、交互题等未明确需求。

\section{Polygon-like 总体业务流程设计}

本项目参考 Polygon 的出题验题流程，但只保留 MVP 所需的本地单题流程。Polygon 中 validator、generator、checker、solutions 和 tests 形成闭环，输出可由标准解自动生成。本项目继承这一思想：generator 只负责生成输入，validator 负责严格校验输入，solution 负责生成输出，系统记录全过程并导出数据包。

\begin{lstlisting}
创建单题项目
  -> 标程导入
  -> 填写数据范围文本
  -> 填写数量规模数值
  -> 选择难度等级下拉选项
  -> Dify 分析标程输入结构与约束草案
  -> 数据结构识别
  -> 测试点计划
  -> 生成 generator.cpp
  -> 生成 validator.cpp
  -> Docker 编译
  -> generator 生成 input
  -> validator 校验 input
  -> main solution 生成 output
  -> 记录资源使用
  -> 生成 report
  -> 导出 zip
\end{lstlisting}

\subsection{流程说明}

\begin{enumerate}
  \item 输入配置阶段以标程代码、数据范围文本、数量规模数值和难度等级下拉选项为入口，不要求完整题面。
  \item 数据范围使用多行文本框填写，形式可接近竞赛题面中的 Constraints，页面必须提示变量名与标程代码保持一致。
  \item 数量规模使用单行文本框填写测试点数量等数值，难度等级使用单选下拉框选择。
  \item AI 分析阶段只产生草案和建议，不能替代确定性校验。
  \item 代码生成阶段输出 generator.cpp 和 validator.cpp，用户可人工检查。
  \item 执行阶段统一进入 Docker 沙箱，禁止宿主机直接运行用户代码。
  \item 验证阶段以 testlib validator 为主，系统基础检查为辅。
  \item 输出阶段只使用用户导入的标程生成 \texttt{.out}。
  \item 报告阶段展示成功结果、失败原因和二次确认问题。
\end{enumerate}

\section{系统架构设计}

\subsection{分层架构}

\begin{longtable}{p{3.2cm}p{10.8cm}}
\toprule
层次 & 职责 \\
\midrule
前端 UI 层 & 项目列表、输入配置、标程导入、资源限制、智能体流程、测试点计划、代码预览、数据预览、报告与导出。 \\
后端 API 层 & 使用 FastAPI 提供项目管理、输入配置管理、标程管理、资源限制、Dify 工作流、代码生成、Docker 执行、数据导出接口。 \\
Dify 智能体工作流层 & 通过 Python 封装 DifyWorkflowClient；优先参考本地 dify 文件夹；不存在可用服务时使用 MockDifyWorkflowClient。 \\
Docker 沙箱层 & 编译和运行 generator、validator、solution；禁止网络；限制 CPU、内存和时间；只挂载临时目录。 \\
testlib 层 & validator.cpp 使用 testlib.h；checker.cpp 暂不实现复杂逻辑，只保留接口。 \\
jngen 层 & generator.cpp 使用 jngen.h；按数据结构选择 array、tree、graph、DAG、matrix、queries 等生成方式。 \\
存储层 & SQLite 保存项目、配置、资源限制和状态；本地 storage 保存 generated files、workdirs、datasets 和 reports。 \\
日志层 & 记录 Dify 调用、Docker 编译、运行、validator 失败和标程运行失败原因。 \\
\bottomrule
\end{longtable}

\subsection{Mermaid 架构图描述}

\begin{lstlisting}[language={}]
flowchart TB
  User[教师 / 出题人] --> UI[React + Vite 前端]
  UI --> API[FastAPI 后端 API]
  API --> DB[(SQLite)]
  API --> FS[本地 storage]
  API --> Dify[DifyWorkflowClient / MockDify]
  API --> Sandbox[Docker 沙箱执行器]
  Dify --> Agents[分析 / 识别 / 规划 / 代码草稿 / 修复 / 报告]
  Sandbox --> Gen[generator.cpp + jngen.h]
  Sandbox --> Val[validator.cpp + testlib.h]
  Sandbox --> Sol[solution.cpp 标程]
  Gen --> Inputs[*.in]
  Inputs --> Val
  Inputs --> Sol
  Sol --> Outputs[*.out]
  Val --> Report[质量报告]
  Outputs --> Package[zip 导出包]
  Report --> Package
\end{lstlisting}

\section{智能体工作流设计}

\subsection{工作流原则}

智能体工作流负责分析与辅助生成，不负责最终事实判定。所有 AI 输出均应进入人工可见的确认环节，并通过 Docker 编译、testlib 校验和标程运行结果进行确定性验证。后端通过 Python 封装 DifyWorkflowClient，优先检查本地 dify 文件夹中的已有调用方式；若不可用，则提供 MockDifyWorkflowClient 保证 MVP 可演示。

环境变量包括：\texttt{DIFY\_BASE\_URL}、\texttt{DIFY\_API\_KEY}、\texttt{DIFY\_WORKFLOW\_ID}、\texttt{USE\_MOCK\_DIFY}。API Key 不允许写死在代码中。

\subsection{节点设计}

\begin{longtable}{p{4.8cm}p{4.5cm}p{4.6cm}}
\toprule
节点 & 输入 & 输出 \\
\midrule
Solution Input Analysis Agent & C++ 标程代码、数据范围、数量规模、难度等级 & 输入结构草案、变量范围草案、结构约束草案、隐含条件、不确定字段 \\
Data Structure Recognition Agent & 输入结构草案、标程读取逻辑 & 需要生成的数据结构，如 array、tree、graph、queries \\
Case Planning Agent & 数据范围、数量规模、难度等级、数据结构识别结果 & 测试点计划，包含 category、scale、purpose、seed、resource limit \\
Generator Code Agent & 测试点计划、数据结构规格 & 基于 jngen 的 generator.cpp 草稿 \\
Validator Code Agent & 输入结构、范围、规模和结构约束 & 基于 testlib 的 validator.cpp 草稿 \\
Repair Agent & 编译错误、运行错误、validator 错误 & 修复后的代码草稿，受最大重试次数限制 \\
Report Agent & 全流程日志、计划、结果、失败原因 & 质量报告和需求方确认报告草稿 \\
\bottomrule
\end{longtable}

\section{Docker 沙箱设计}

\subsection{镜像内容}

Docker 镜像应包含 Linux 基础环境、g++、make、Python 运行时、testlib.h、jngen.h 以及必要的脚本工具。镜像不包含任何业务密钥，不在容器内访问外部网络。

\subsection{编译流程}

\begin{enumerate}
  \item 后端为每次任务创建临时 workdir。
  \item 写入 generator.cpp、validator.cpp、solution.cpp 和资源头文件。
  \item 使用 Docker 运行编译命令，分别生成 generator、validator、solution 可执行文件。
  \item 捕获 stdout、stderr、exit code 和编译耗时。
  \item 任一核心代码编译失败时停止后续数据生成。
\end{enumerate}

\subsection{运行流程}

\begin{enumerate}
  \item 按测试点计划调用 generator，传入 seed、case\_id、scale 等参数。
  \item 对生成的 \texttt{.in} 执行文件大小检查和基础格式检查。
  \item 使用 validator 校验输入文件。
  \item 使用 solution 读取 \texttt{.in} 并生成 \texttt{.out}。
  \item 记录运行时间、内存、退出码、stdout size、stderr、是否超时、是否内存超限。
\end{enumerate}

\subsection{安全限制}

\begin{itemize}
  \item 网络禁用：Docker 运行使用 \texttt{--network none}。
  \item 内存限制：通过 \texttt{--memory} 设置容器上限。
  \item CPU 限制：通过 \texttt{--cpus} 设置可用 CPU。
  \item 文件挂载：只挂载临时工作目录，不挂载项目根目录和用户目录。
  \item 时间限制：后端对子进程设置 timeout，超时后终止容器。
\end{itemize}

\subsection{错误分类}

\begin{longtable}{p{5.2cm}p{8.8cm}}
\toprule
错误码 & 含义 \\
\midrule
\texttt{COMPILE\_ERROR} & generator、validator 或 solution 编译失败。 \\
\texttt{RUNTIME\_ERROR} & 程序运行时异常退出。 \\
\texttt{TIME\_LIMIT\_EXCEEDED} & 运行超过配置时间限制。 \\
\texttt{MEMORY\_LIMIT\_EXCEEDED} & 运行超过配置内存限制。 \\
\texttt{VALIDATION\_FAILED} & 输入文件未通过 testlib validator。 \\
\texttt{INPUT\_TOO\_LARGE} & 输入文件超过单文件大小限制。 \\
\texttt{OUTPUT\_TOO\_LARGE} & 输出文件超过单文件大小限制。 \\
\texttt{DOCKER\_ERROR} & Docker 启动、挂载或资源控制异常。 \\
\bottomrule
\end{longtable}

\section{内存控制设计}

\subsection{资源限制字段}

\begin{longtable}{p{5.8cm}p{8.4cm}}
\toprule
字段 & 说明 \\
\midrule
\texttt{generator\_memory\_mb} & generator 运行内存上限。 \\
\texttt{validator\_memory\_mb} & validator 运行内存上限。 \\
\texttt{solution\_memory\_mb} & 标程运行内存上限。 \\
\texttt{generator\_timeout\_ms} & generator 单次运行超时。 \\
\texttt{validator\_timeout\_ms} & validator 单次运行超时。 \\
\texttt{solution\_timeout\_ms} & 标程单个测试点运行超时。 \\
\texttt{max\_input\_file\_size\_mb} & 单个输入文件最大大小。 \\
\texttt{max\_output\_file\_size\_mb} & 单个输出文件最大大小。 \\
\texttt{max\_total\_dataset\_size\_mb} & 最终数据包总大小上限。 \\
\texttt{max\_retry\_count} & AI 修复和重新生成的最大重试次数。 \\
\bottomrule
\end{longtable}

\subsection{失败停止策略}

后端任务采用 fail-fast 策略。标程编译失败、validator 编译失败、generator 编译失败时终止流程。单个测试点生成失败时可按最大重试次数重试；重试后仍失败则标记为 invalid，不进入最终数据包。若合法测试点数量低于用户要求，则报告中明确提示“数据包不满足数量规模要求”。

\section{jngen 数据结构生成设计}

本项目不按“数组题、图题、字符串题”等题型分类，而按“需要生成的数据结构”分类。jngen README 显示其支持随机数、命令行参数、数组包装、打印器、图和树生成、字符串、几何对象等能力，适合作为复杂结构 generator 的基础库。

\begin{longtable}{p{2.5cm}p{3.4cm}p{4.2cm}p{4cm}}
\toprule
结构 & 适用输入 & 生成策略 & validator 校验点 \\
\midrule
array & 一维数组、权值列表 & 随机、单调、重复值、极值混合 & 长度、元素范围、换行和空格 \\
sequence & 有序序列、操作序列 & 递增、递减、震荡、局部有序 & 顺序关系、长度、范围 \\
permutation & 排列、排名、映射 & 随机排列、逆序、近似有序 & 元素唯一性、取值为 1..n \\
string & 字符串、模式串 & 随机字符、重复块、边界长度 & 字符集、长度、格式 \\
tree & 树结构 & 链、星、随机树、重心偏移树 & n-1 条边、连通、无环 \\
graph & 一般图 & 稀疏、稠密、连通、多连通块 & 点边范围、重边、自环策略 \\
DAG & 有向无环图 & 拓扑序生成、分层 DAG & 无环、边方向、点边范围 \\
bipartite graph & 二分图、匹配图 & 两侧规模可控、随机边、极端不均衡 & 左右集合边界、边合法性 \\
matrix & 矩阵、二维权值 & 随机、行列模式、极值块 & 行列数、元素范围 \\
grid & 网格地图 & 障碍密度、连通区域、边界墙 & 行列数、字符集、连通性要求 \\
queries & 查询序列 & 更新/询问混合、极端重复、在线顺序 & 查询数、操作类型、参数范围 \\
geometry points & 平面点集 & 随机点、共线、凸包、极值坐标 & 坐标范围、重复点策略 \\
combined structure & 图+查询、树+权值等 & 分阶段生成并保持关联关系 & 各子结构约束与跨字段关系 \\
\bottomrule
\end{longtable}

\section{testlib validator 设计}

testlib README 显示，validator 通过 \texttt{registerValidation(argc, argv)} 读取标准输入并严格检查格式、范围、行尾和 EOF；非法输入会返回非零状态并输出错误信息。MVP 中 validator 是主要验题标准，系统基础校验器只做文件大小、空文件、基本可读性等辅助检查。

\subsection{validator 职责}

\begin{enumerate}
  \item 检查输入格式是否与确认后的输入结构一致。
  \item 检查整数、字符串、图边、树边、查询参数等是否在数据范围内。
  \item 检查结构约束，如树连通无环、排列唯一、图边合法、DAG 无环。
  \item 检查输入结尾，避免多余 token。
  \item 输出明确错误提示，便于 RepairAgent 或人工修改 generator。
\end{enumerate}

\subsection{与基础校验器的关系}

系统基础校验器负责轻量检查：文件是否存在、是否为空、是否超过大小限制、是否重复、是否包含不可接受的二进制内容。testlib validator 负责领域语义校验，是判断 \texttt{.in} 是否可进入最终数据包的核心依据。

\section{数据模型设计}

\begin{longtable}{p{4.2cm}p{9.8cm}}
\toprule
模型 & 关键字段与说明 \\
\midrule
Project & project\_id、name、description、status、created\_at、updated\_at。 \\
ProblemInputConfig & standard\_code\_summary、range\_spec、scale\_spec、difficulty、confirmed\_fields、uncertain\_fields。 \\
ResourceLimit & generator/validator/solution 的 memory、timeout、CPU、文件大小和总包大小限制。 \\
StandardSolution & source\_path、language、compile\_status、compile\_log、binary\_path。 \\
DataStructureSpec & structure\_type、variables、constraints、confidence、needs\_confirmation。 \\
CasePlan & case\_id、category、scale、purpose、seed、resource\_limit。 \\
GeneratedCode & code\_type、source\_path、compile\_status、repair\_count、latest\_log。 \\
TestCase & case\_id、input\_path、output\_path、category、status、hash。 \\
ValidationResult & case\_id、basic\_check、validator\_exit\_code、message、passed。 \\
ExecutionResult & case\_id、exit\_code、runtime\_ms、memory\_kb、stdout\_size、stderr、verdict。 \\
ResourceUsage & case\_id、phase、time\_ms、memory\_kb、file\_size\_bytes。 \\
WorkflowRunLog & run\_id、node\_name、input\_summary、output\_summary、status、error。 \\
ExportPackage & package\_id、zip\_path、metadata\_path、report\_path、created\_at。 \\
\bottomrule
\end{longtable}

\section{API 设计}

\begin{longtable}{p{6.8cm}p{7.2cm}}
\toprule
接口 & 说明 \\
\midrule
\footnotesize\url{POST /api/projects} & 创建项目。 \\
\footnotesize\url{GET /api/projects} & 获取项目列表。 \\
\footnotesize\url{GET /api/projects/{project_id}} & 获取项目详情。 \\
\footnotesize\url{PUT /api/projects/{project_id}} & 更新项目基础信息。 \\
\footnotesize\url{POST /api/projects/{project_id}/solution} & 上传或保存 C++ 标程。 \\
\footnotesize\url{POST /api/projects/{project_id}/solution/compile} & Docker 编译标程。 \\
\footnotesize\url{PUT /api/projects/{project_id}/input-config} & 保存数据范围、数量规模、难度等级。 \\
\footnotesize\url{GET /api/projects/{project_id}/input-config} & 查询输入配置。 \\
\footnotesize\url{PUT /api/projects/{project_id}/resource-limits} & 保存资源限制。 \\
\footnotesize\url{GET /api/projects/{project_id}/resource-limits} & 查询资源限制。 \\
\footnotesize\url{POST /api/projects/{project_id}/workflow/analyze} & 运行 Dify / Mock Dify 分析流程。 \\
\footnotesize\url{GET /api/projects/{project_id}/workflow/result} & 获取工作流结果。 \\
\footnotesize\url{POST /api/projects/{project_id}/generate} & 生成并校验数据。 \\
\footnotesize\url{GET /api/projects/{project_id}/cases} & 获取测试点列表。 \\
\footnotesize\url{GET /api/projects/{project_id}/cases/{case_id}} & 获取单个测试点详情。 \\
\footnotesize\url{GET /api/projects/{project_id}/report} & 获取质量报告。 \\
\footnotesize\url{POST /api/projects/{project_id}/export} & 导出 zip 数据包。 \\
\bottomrule
\end{longtable}

\section{界面设计}

\begin{longtable}{p{3.2cm}p{4.5cm}p{6.3cm}}
\toprule
页面 & 页面目标 & 核心组件与交互 \\
\midrule
首页 / 项目列表页 & 管理单题项目 & 项目表格、状态标签、新建按钮、搜索筛选、最近运行状态。 \\
输入配置页 & 汇总本题输入入口 & 标程状态、数据范围文本框、数量规模文本框、难度等级单选下拉框、确认状态。 \\
标程导入页 & 提供唯一答案来源 & 代码编辑器、上传按钮、编译按钮、编译日志、错误定位。 \\
数据范围配置页 & 录入变量范围和结构关系 & 多行文本框、Constraints 示例提示、变量名一致性提示。 \\
数量规模配置页 & 配置测试点数量等数值规模 & 单行文本框、数字格式校验、默认值提示。 \\
难度等级配置页 & 控制测试强度 & 单选下拉框、难度说明、覆盖强度预览。 \\
资源限制配置页 & 控制时间、内存、文件大小 & timeout、memory、CPU、输入输出大小、总包大小。 \\
智能体流程页 & 展示 Dify 节点执行状态 & 节点时间线、输入输出摘要、不确定字段、重试记录。 \\
测试点计划页 & 人工确认测试点覆盖 & category、scale、purpose、seed 表格，启用/禁用开关。 \\
代码预览页 & 查看 generator / validator & 双栏代码预览、编译状态、修复建议、人工确认按钮。 \\
数据预览页 & 查看输入输出和日志 & 测试点列表、输入预览、输出预览、资源使用。 \\
质量报告页 & 展示最终质量结论 & 覆盖摘要、失败原因、人工确认事项、导出入口。 \\
导出页 & 生成本地数据包 & zip 内容清单、导出按钮、下载路径、metadata 预览。 \\
\bottomrule
\end{longtable}

\subsection{前端输入控件设计}

输入配置相关控件应尽量降低用户填写成本，同时保证后端能够稳定解析。数据范围不采用复杂表单拆分，而采用多行文本框，让用户输入接近竞赛题面 Constraints 的文本描述。页面需要给出示例，例如 \texttt{2 <= N <= 5 * 10\^{}5}、\texttt{N is an integer}、\texttt{S is a string of length N consisting of o and x}。输入框旁必须提示：变量名需要与标程代码中的变量名保持一致，否则系统分析标程读取逻辑时可能无法对应字段。

数量规模采用单行文本框，用户输入测试点数量等数值规模，前端需要做非空、正整数和合理上限校验。难度等级采用单选下拉框，建议提供“简单”“中等”“较难”三个选项，用户每次只能选择一个难度等级。

\begin{longtable}{p{3cm}p{3.8cm}p{6.8cm}}
\toprule
配置项 & 控件形式 & 交互与校验要求 \\
\midrule
数据范围 & 多行文本框 & 支持类似 Constraints 的自然语言或约束列表；提示变量名必须与标程代码一致；用于 AI 分析输入结构与 validator 草稿。 \\
数量规模 & 单行文本框 & 用户输入测试点数量等数值；前端校验正整数、非空值和上限；后端据此生成测试点计划。 \\
难度等级 & 单选下拉框 & 从预设难度中选择一个；影响边界数据、极限数据和特殊构造比例。 \\
\bottomrule
\end{longtable}

\section{应用功能边界}

\subsection{建议纳入应用的功能}

\begin{longtable}{p{4cm}p{9.5cm}}
\toprule
功能 & 说明 \\
\midrule
单题项目管理 & 创建、查看、编辑单题项目，保存项目状态和基础配置。 \\
标程导入与编译 & 上传 C++ 标程并在 Docker 中检查是否能编译。 \\
数据范围配置 & 通过多行文本框录入类似 Constraints 的数据范围说明，变量名需与标程代码保持一致。 \\
数量规模配置 & 通过单行文本框录入测试点数量等数值规模，并做数字格式校验。 \\
难度等级配置 & 通过单选下拉框选择难度等级。 \\
AI 分析与规划 & 分析标程输入结构，生成输入结构草案、数据结构识别结果和测试点计划。 \\
数据结构识别 & 按 array、tree、graph、queries 等数据结构识别，而不是按题型识别。 \\
输入生成 & 使用 generator 批量生成 \texttt{.in}。 \\
输入校验 & 使用 validator 检查格式、范围和逻辑合理性。 \\
标准输出生成 & 使用用户导入的标程生成 \texttt{.out}。 \\
质量报告 & 展示覆盖情况、失败原因和资源使用。 \\
zip 导出 & 导出 \texttt{.in/.out}、代码和报告。 \\
\bottomrule
\end{longtable}

\subsection{暂不纳入应用需要二次确认的功能}

\begin{longtable}{p{4cm}p{9.5cm}}
\toprule
功能或事项 & 暂不纳入原因或需确认问题 \\
\midrule
OJ 平台兼容 & 无法理解需求方的“兼容 OJ 平台、课堂练习、模拟比赛”是什么具体需求，需要确认兼容对象是文件命名、zip 结构、metadata、checker / validator，还是具体 OJ 的 API 对接。 \\
课堂练习 / 模拟比赛 & 暂不纳入当前应用。需确认需求方是否需要教师下载数据、学生在线练习、多题管理、分值、排名和提交记录等完整教学或比赛流程。 \\
SPJ 自动生成 & 暂不纳入当前应用。特殊评测通常需要完整题面和输出判定规则，仅通过标程代码难以可靠推断。 \\
交互题 & 暂不纳入当前应用。交互题的执行、输入输出协议和校验机制与普通题不同，需要单独设计。 \\
完整题面辅助输入 & 当前应用默认不输入完整题面。需确认后续是否允许把题面作为可选辅助信息，以提升 AI 对输入结构和特殊规则的理解。 \\
文件与测试点上限 & 需确认单题最大测试点数量、单个输入输出文件大小和总数据包大小上限。 \\
\bottomrule
\end{longtable}

\section{验收标准}

\begin{longtable}{p{1.2cm}p{12.6cm}}
\toprule
序号 & 可检查标准 \\
\midrule
1 & 能创建单题项目。 \\
2 & 能导入并编译 C++ 标程。 \\
3 & 能录入数据范围、数量规模、难度等级。 \\
4 & 能分析标程输入结构并输出需要人工确认的不确定字段。 \\
5 & 能生成测试点计划。 \\
6 & 能生成 generator.cpp。 \\
7 & 能生成 validator.cpp。 \\
8 & generator 和 validator 都在 Docker 中运行。 \\
9 & 每个 \texttt{.in} 通过 validator。 \\
10 & 标程能生成 \texttt{.out}。 \\
11 & 能显示时间、内存、文件大小。 \\
12 & 能导出 zip。 \\
13 & 能生成 report.md。 \\
14 & 遇到编译错误、运行错误、超时、内存超限时能显示明确错误。 \\
\bottomrule
\end{longtable}

\section{推荐技术栈}

\begin{itemize}
  \item 后端：FastAPI。
  \item 前端：React + Vite。
  \item 数据库：SQLite。
  \item 执行环境：Docker。
  \item 智能体调用：Dify Python client wrapper。
  \item 输入校验：testlib.h。
  \item 数据生成：jngen.h。
  \item 文件存储：本地 storage。
  \item 导出：Python zipfile。
  \item Docker 调用：后端 subprocess 调 Docker，但不直接运行用户代码。
\end{itemize}

\section{项目目录结构建议}

\begin{lstlisting}
backend/
  app/
    api/
    models/
    services/
    sandbox/
    workflows/
frontend/
  src/
docker/
third_party/testlib/
third_party/jngen/
dify/
storage/
  projects/
  workdirs/
  datasets/
  reports/
docs/
scripts/
\end{lstlisting}

\section{小组分工}

\begin{longtable}{p{3.5cm}p{10.5cm}}
\toprule
方向 & 工作内容 \\
\midrule
前端 UI & 页面设计、交互流程、代码预览、数据预览、报告展示。 \\
后端 API & FastAPI 接口、数据模型、任务状态、导出逻辑。 \\
Docker 沙箱 & 镜像、编译运行、资源限制、错误分类。 \\
Dify 工作流 & Client 封装、Mock 实现、节点输入输出、日志记录。 \\
jngen generator & 数据结构生成策略、seed/case\_id、生成参数。 \\
testlib validator & 输入格式、范围、结构约束、错误提示。 \\
报告与答辩 & 设计报告、需求确认报告、演示脚本、PPT 素材。 \\
\bottomrule
\end{longtable}

\section{开发里程碑}

\begin{longtable}{p{2.2cm}p{11.8cm}}
\toprule
时间 & 任务 \\
\midrule
Day 1 & 需求确认、架构、报告初稿。 \\
Day 2 & 前后端骨架、数据库模型。 \\
Day 3 & Docker 沙箱、标程编译运行。 \\
Day 4 & jngen generator、testlib validator。 \\
Day 5 & Dify / Mock Dify 工作流。 \\
Day 6 & 数据生成、校验、报告、导出。 \\
Day 7 & 联调、界面优化、演示准备。 \\
\bottomrule
\end{longtable}

\section{参考资料与当前目录依据}

\begin{itemize}
  \item \texttt{goal.md}：定义项目定位、MVP 边界、流程、架构、API、界面和验收要求。
  \item \texttt{课程考查要求.pdf}：明确项目设计报告需要覆盖业务逻辑、功能设计、交互界面和二次需求确认。
  \item \texttt{项目要求.jpg}：明确项目题目和主要功能，包括标程代码、数据范围、数量规模、难度等级、数据生成、数据校验和导出。
  \item \texttt{testlib/README.md}：说明 testlib 在 Codeforces 等竞赛中的使用，并提供 validator、checker、generator 的典型模式。
  \item \texttt{jngen/README.md}：说明 jngen 支持随机数、参数解析、数组、图、树、字符串和几何等生成能力。
  \item \texttt{dify/README.md}：说明 Dify 是开放的 LLM 应用开发平台，支持 Workflow、Agent、API 和 Docker Compose 本地部署。
  \item \texttt{OI-wiki/docs/tools/polygon.md}：说明 Polygon 中 validator、generator、checker、solutions、Tests、Invocations、Packages 等出题验题流程。
\end{itemize}

\end{document}
